\section{Local spin and momentum-dependent visibility}
\label{sec:results}
\label{sec:doublet}

\subsection{Projection onto a cell spin}

The isolated cell retains a half-integer spin for either sign of the outer
bonds. All 100 sampled points with $0<\alpha\le1$ have an even-parity
spin-$1/2$ ground doublet; all 201 sampled points with
$-1\le\alpha\le-0.001$ have a spin-$3/2$ ground quartet.
The different local spins provide two realizations of an antiferromagnetically
coupled cluster chain, with different internal spin distributions.

Let $P$ project onto one isolated ground multiplet of spin $s$.
Within this subspace rotational symmetry gives
\begin{equation}
 P\bm S_{n,a}P=c_a\bm T_n,\qquad \sum_{a=1}^5c_a=1,
 \label{eq:projection}
\end{equation}
where $\bm T_n$ is a spin-$s$ operator.
The coefficients $c_a$ specify where the retained spin resides within a cell;
negative coefficients indicate a local moment antiparallel to the cell spin.
Reflection implies $c_a=c_{6-a}$ for either reflection parity, because the
parity phases of the bra and ket cancel.
Consequently the first-order Hamiltonian is
\begin{align}
 \Heff&=NE_5+\Jeff\sum_n\bm T_n\!\cdot\!\bm T_{n+1},\notag\\
 \Jeff&=\beta c_1c_5=\beta c_1^2>0.
 \label{eq:effective}
\end{align}
Here $E_5$ is the cell ground-state energy; corrections are of order
$\beta^2/\Delta_{\mathrm{loc}}$ at fixed cell parameters, where
$\Delta_{\mathrm{loc}}$ is the gap above the entire ground multiplet.
The controlled limit is $g=\beta/\Delta_{\mathrm{loc}}\ll1$.

Projecting the probe, rather than only the Hamiltonian, reveals the observable
consequence. Its intra-cell amplitude is
$f(q)=\sum_ac_ae^{iqa}$ and its shape factor is
\begin{equation}
 F(q)=|f(q)|^2.
 \label{eq:shape}
\end{equation}
With the physical and effective probes normalized by $1/\sqrt L$ and
$1/\sqrt N$, respectively,
\begin{equation}
 S^{zz}_{\mathrm{phys}}(q,\omega)\simeq
 \frac{F(q)}5 S^{zz}_{\mathrm{eff}}(Q,\omega),
 \quad Q=5q\pmod{2\pi}.
 \label{eq:folded-response}
\end{equation}
The factor $1/5$ is fixed by normalization, not fitted.
Thus replicas at the same $Q$ have a common collective response but unequal
visibility set by $F(q)$.

Table~\ref{tab:projection} makes this contrast quantitative at
$\alpha=\pm0.5$.
For positive $\alpha$ the $q=\pi$ replica is brightest; for negative $\alpha$
it is suppressed by almost three orders of magnitude relative to $0.60\pi$.
The latter is the largest of the three folded antiferromagnetic replicas,
not the global maximum of $F(q)$: $F(0)=(\sum_ac_a)^2=1$.
The location of a peak in the full spectrum depends on the product in
Eq.~\eqref{eq:folded-response}, not on $F$ alone.

\begin{table}[!htbp]
 \caption{Isolated-cell projection at the two representative couplings.
 Reflection fixes $c_5=c_1$ and $c_4=c_2$. The last three rows compare
 physical momenta with the same cell momentum $Q=\pi$.}
 \label{tab:projection}
 \begin{ruledtabular}
 \begin{tabular}{lrr}
 & $\alpha=0.5$ & $\alpha=-0.5$\\
 $s$ & $1/2$ & $3/2$\\
 $\Delta_{\mathrm{loc}}/J_1$ & 0.342371 & 0.126397\\
 $c_1$ & 0.622839 & 0.331329\\
 $c_2$ & $-0.277350$ & 0.242242\\
 $c_3$ & 0.309022 & $-0.147142$\\
 $\Jeff/\beta$ & 0.387928 & 0.109779\\
 $F(0.20\pi)$ & 0.060121 & 0.202128\\
 $F(0.60\pi)$ & 0.278088 & 0.693819\\
 $F(\pi)$ & 4.449570 & 0.000963\\
 \end{tabular}
 \end{ruledtabular}
\end{table}

\subsection{An absolute two-cell spectral-weight test}
\label{sec:projection-benchmark}

A two-cell chain provides a direct test of both the projected energy scale
and the probe normalization, without assumptions about a thermodynamic line
shape. Two antiferromagnetically coupled spins $s$ form a singlet, with
$\langle(T_1^z)^2\rangle=s(s+1)/3$ and
$\langle T_1^zT_2^z\rangle=-s(s+1)/3$.
Their projected physical probe is
$f(q)[T_1^z+e^{i5q}T_2^z]/\sqrt{10}$.
Its norm and the singlet--triplet splitting therefore predict
\begin{align}
 W_{\mathrm{low}}^{(1)}(q)
 &=\frac{F(q)s(s+1)}{15}[1-\cos(5q)],
 \label{eq:two-cell-weight}\\
 \Delta_{10}^{(1)}&=\beta c_1c_5.
 \label{eq:two-cell-gap}
\end{align}
There is no adjustable amplitude or independent normalization at each momentum.

We compare these expressions with complete exact diagonalization in the
$S^z=0$ sector (dimension 252) of two full pentamers and the lowest energy in
$S^z=1$ (dimension 210).
Before diagonalization, the low band is defined as
$0<\omega<\Delta_{\mathrm{loc}}/2$.
The calculation uses $g=0.0025,0.005,0.01,0.02,0.04$ at each sign of $\alpha$;
$g\le0.01$ is the predefined weak domain.
Both the spectral completeness sum rule and the Fourier-averaged local moment
are checked to absolute error below $10^{-11}$.

Figure~\ref{fig:projection-benchmark} shows that Eq.~\eqref{eq:two-cell-weight}
reproduces the dominant momentum pattern and its suppressed replicas.
Across the six weak-domain cases the largest gap error is
$0.5574\%$ and the largest relative $L_2$ error of the low-weight vector is
$0.2643\%$, below the predefined $2\%$ and $3\%$ tolerances.
At $g=0.01$, the positive and negative weight errors are $0.2643\%$ and
$0.0884\%$, respectively.
The vector error is not a uniform pointwise bound: relative errors can be
larger where destructive interference makes the predicted intensity tiny.
This test verifies absolute low-band weights as well as the effective energy
scale; the reversal of replica brightness is not a consequence of independently
rescaling the spectra.

\begin{figure*}[t]
 \centering
 \includegraphics[width=0.96\textwidth]{fig_projection_ed_benchmark.pdf}
 \caption{Absolute low-energy projection test for two pentamers ($L=10$).
 (a,b) Low-band connected weights at $g=0.01$ for $\alpha=0.5$ and $-0.5$.
 Circles are exact-diagonalization values at 41 momenta; lines join
 Eq.~\eqref{eq:two-cell-weight} at the same momenta, without a fitted amplitude.
 The band $0<\omega<\Delta_{\mathrm{loc}}/2$ is fixed before the calculation.
 (c,d) Gap error $|\Delta_{\mathrm{ED}}/\Delta_{10}^{(1)}-1|$ and weight error
 $\|\mathbf W_{\mathrm{ED}}-\mathbf W^{(1)}\|_2/\|\mathbf W^{(1)}\|_2$.
 Shading marks the predefined domain $g\le0.01$; segments connect only the
 five sampled couplings. This small-system test does not establish
 thermodynamic convergence. All symbols
 and lines are rendered from the calculated data and analytical prediction.}
 \label{fig:projection-benchmark}
\end{figure*}

\subsection{A distinct parity filter for local excitations}
\label{sec:parity}

The low-band interference above and the selection of excited local multiplets
are different effects.
At $q=\pi$, the alternating signs within a five-site cell are invariant under
$a\mapsto6-a$. The local spin probe is therefore reflection-even and cannot
connect states of opposite parity.
For positive $\alpha$ the ground doublet is even and the lowest excited
doublet is odd. This excited doublet is dark, although it is the lowest local
excitation in energy.
At $\alpha=0.14,0.23,0.33$, its multiplet-summed weight is below $10^{-25}$,
whereas spin-$3/2$ multiplets carry $84\%$--$98\%$ of the local inelastic
weight (Appendix~\ref{sec:local-details}).
Reflection selects which local states can be reached; it does not alone
identify a dispersive band after the cells are strongly coupled.

\section{Finite-coupling spectra and momentum reconstruction}
\label{sec:dynamics}

\subsection{Full-momentum response at fixed resolution}

Figure~\ref{fig:fullq} compares the connected spectra at
$(\alpha,\beta)=(0.5,0.5)$ and $(-0.5,0.5)$ for $L=40$.
At positive $\alpha$, low-energy weight is prominent near $q=\pi$ and coexists
with higher-energy bands. At negative $\alpha$, the low-energy response shifts
toward other momenta and the higher-energy pattern also changes.
All 41 momentum sources are independently calculated, and their reflection
deviations are only $2.87\times10^{-6}$ and $4.49\times10^{-6}$ in relative
root-mean-square norm. The symmetry is thus a numerical check, not imposed
by averaging the left and right halves.

These spectra extend the comparison to substantial inter-cell coupling, but
do not extend the quantitative validity of the weak-coupling formula.
Here $g=1.4604$ and $3.9558$ for the two signs.
Two-cell diagonalization at the same $\beta=0.5$ fails the predefined
low-band separation criterion for both signs; the first-order gap errors
increase to approximately $10\%$ and $84\%$.
The weak-limit replica weights therefore provide a reference for reading the
redistribution, not a quantitative fit or a unique assignment of these bands.

The energy resolution is $0.09277J_1$.
Independent sector calculations give open-chain gaps $0.05562J_1$ and
$0.001092J_1$, both narrower than this width.
The maps can compare weight distributions on the displayed energy scale,
but do not establish that either gap is resolved from zero frequency.
Weak stripes enhanced by the common color scale must likewise not be counted
as additional excitations on their appearance alone.

\begin{figure*}[t]
 \centering
 \includegraphics[width=0.49\textwidth]{fig6_fullq_positive.pdf}\hfill
 \includegraphics[width=0.49\textwidth]{fig6_fullq_negative.pdf}
 \caption{Complete connected $S_T^{zz}(q,\omega)$ for $L=40$ open chains at
 $J_2/J_1=0.5$ (left) and $-0.5$ (right), with $J_3/J_1=0.5$.
 The physical momentum spans $0\le q\le2\pi$ and energies are in units of $J_1$.
 The triangular time window at $TJ_1=60$ has width $0.09277J_1$;
 production settings are $\Delta tJ_1=0.05$ and $\chi=192$.
 Both panels use the same blue-to-red color coordinate
 $u(s)=\operatorname{asinh}(s/0.02)/\operatorname{asinh}(6/0.02)$ for
 $s=J_1S_T^{zz}\in[0,6]$.
 This uniform nonlinear display enhances weak weight without changing the
 spectral values or their energy resolution. No Gaussian smoothing,
 momentum interpolation, or reflection averaging is used.
 The panels render the calculated
 spectral arrays.}
 \label{fig:fullq}
\end{figure*}

\subsection{Finite-size reconstruction on the ferromagnetic side}
\label{sec:negative}

The static data ask a complementary question: as $\alpha$ approaches zero from
below at fixed $\beta$, does the dominant response retain the $0.60\pi$
momentum favored among the weak-limit antiferromagnetic replicas?
The negative-$\alpha$ scan instead finds a movement toward $q=\pi$.
Figure~\ref{fig:negative-reconstruction} displays the converged $L=80$
samples and their $L=60$ comparison, without interpolating a phase surface.

To test whether this movement locates a common boundary, we require an adjacent
coupling interval to show both a change in $\qpeak$ and a sign change of
$\Delta R$, with $|\Delta R|\ge2\epsilon$ at both endpoints.
Here $\epsilon$ is a deterministic numerical tolerance, built from the
initial-state spread, truncation error, and a $10^{-4}$ floor; it is not a
statistical standard deviation.
Of eight candidate intervals across seven $\beta$ slices, seven retain at
least one signature, but none has both. Only the $\beta=0.05$ momentum
interval reaches the target width $0.025$.
The $\beta=0.10$ slice supplies no supported internal interval at the final
size pair.

The dominant momentum is therefore changing in these finite chains, while the
two boundary diagnostics do not coincide.
This establishes a finite-size reconstruction rather than a separate
thermodynamic phase. A broad crossover, slow convergence of an incommensurate
regime, and residual finite-size effects remain possible interpretations.
An apparent entropy suppression near $(\alpha,\beta)=(-0.10,0.25)$ also fails
an eigensolver control: the fitted coefficient changes from about $0.53$ to
$0.93$ at $L=80$ (Appendix~\ref{sec:krylov}).
It therefore supplies no independent evidence for a new low-entropy phase.

\begin{figure*}[t]
 \centering
 \includegraphics[width=0.96\textwidth]{fig8_finite_size_reconstruction.pdf}
 \caption{Finite-size momentum reconstruction for negative $\alpha$.
 (a) Dominant momentum at the 22 sampled $L=80$ couplings; crosses mark
 disagreement of the momentum classification between $L=60$ and 80.
 Gray segments connect sampled points only.
 (b) Peak-sharpness flow $\Delta R=R_{80}-R_{60}$.
 Bars show $\epsilon=(\epsilon_{80}^2+\epsilon_{60}^2)^{1/2}$, where each
 size-specific tolerance is the larger of the initial-state half-range,
 ten times the maximum final truncation error, and $10^{-4}$.
 These are not statistical confidence intervals.
 (c) Eight candidate intervals tested for a momentum change, a sign change
 satisfying $|\Delta R|\ge2\epsilon$ at both endpoints, width at most $0.025$,
 and the joint momentum/size-flow criterion. No target satisfies the joint
 criterion; this operational test neither establishes nor rules out a
 thermodynamic transition. The symbols
 and decisions are taken from the recorded numerical analysis.}
 \label{fig:negative-reconstruction}
\end{figure*}
